\scanpage{314}
\mspage{13}
% archive sheet number: 315

donc
\[
\begin{aligned}
\widehat\Delta^{\circ}\xrightarrow{F_0}B_0
=(\widehat\Delta_{\mathrm{ab}})^{\circ}
&\simeq\underline{\operatorname{Hom}}(\Delta^{\circ},\mathit{Ab})^{\circ}\\
&\approx\underline{\operatorname{Hom}}(\Delta,\mathit{Ab}^{\circ})\\
&[\simeq(\mathit{Ab}^{\circ})^{\ast}\inserted{\simeq(\mathit{Ab}_{\ast})^{\circ}},
\end{aligned}
\]
\inserted{dans le cas qui nous intéresse]} commutant aux $\varprojlim$
d'où le foncteur
\[
\underline{\operatorname{Hom}}_{!}(\deleted{\unsure{\widehat\Delta_{\mathrm{ab}}}}B_0,B)
\longrightarrow\underline{\operatorname{Hom}}_{!}(\widehat\Delta^{\circ},B)
\]

\deleted{et on définit un foncteur}
dont on veut montrer que c'est une équiv. Posant
$B=A^{\circ}$, $B_0=(A_0)^{\circ}$, ou $F_0\deleted{\unsure{\longleftrightarrow}}=(G_0)^{\circ}$, donc
\[
G_0:\widehat\Delta\longrightarrow\widehat\Delta_{\mathrm{ab}},\quad G_0(X)=\mathbb Z(X),
\]
on doit prouver que
\[
\begin{aligned}
\underline{\operatorname{Hom}}_{!}^{\unsure{\S|!|\$}}(\widehat\Delta_{\mathrm{ab}},A)&\longrightarrow
\underline{\operatorname{Hom}}_{!}^{\S}(\widehat\Delta,A)\\
\unsure{\varphi|\psi}&\longmapsto\psi\circ G_0
\end{aligned}
\]
est une équiv., où $ {}_{!}^{\S}$ désigne les foncteurs
qui commutent aux $\varinjlim$. Or
le deuxième mb est $\approx\underline{\operatorname{Hom}}(\Delta,A)$
(par restriction à $\Delta\hookrightarrow\widehat\Delta$). \deleted{\unsure{donc}} Le
premier s'écrit ? \deleted{\unsure{en|on}} Prenant $\varphi\in\underline{\operatorname{Hom}}_{!}(\widehat\Delta_{\mathrm{ab}},A)$
et notant que pour tout $M\in A$, \deleted{\unsure{l'applic.}} le
foncteur $X\longmapsto\operatorname{Hom}(\varphi X,\deleted{\unsure{\varphi}}M)$ transforme $\varinjlim$
en $\varprojlim$ donc est représentable
\[
\operatorname{Hom}(\varphi X,M)\simeq\operatorname{Hom}(X,\psi M)
\]

i.e. l'hyp. de commutation de $\varphi$ aux $\varinjlim$
s'exprime par le fait qu'il admet un
adjoint $\psi$ -- la condition à mettre sur \deleted{$\unsure{\varphi}$}
$\psi:A\longrightarrow\widehat\Delta_{\mathrm{ab}}$ pour qu'il provienne de $\varphi$
étant que pour tt $X$ fixé ds $\widehat\Delta_{\mathrm{ab}}$, $M\longrightarrow\operatorname{Hom}(X,\psi M)$
soit repr. Ecrivant $\widehat\Delta_{\mathrm{ab}}\simeq\operatorname{Hom}(\Delta^{\circ},(\mathit{Ab}))$
on trouve donc interprétation de la

cat. $\underline{\operatorname{Hom}}_{!}(\widehat\Delta_{\mathrm{ab}},A)$ des $\varphi$ comme ss-cat. pleine
de $\underline{\operatorname{Hom}}(A\times\Delta^{\circ},\mathit{Ab})$\inserted{$\simeq\underline{\operatorname{Hom}}(\Delta^{\circ},\underline{\operatorname{Hom}}(A,\mathit{Ab}))$}, savoir
des foncteurs $\lambda$, satisfaisant aux conditions de
repr. \inserted{ci-dessus} en $A$. Or si $x\in\Delta$, le foncteur \deleted{$M\longmapsto$}
$\deleted{\lambda(\unsure{X|x})(\unsure{M})}\inserted{\lambda(x):A\longrightarrow}\in\mathit{Ab}$ n'est autre que
\[
M\longmapsto\operatorname{Hom}(\mathbb Z(x),\psi(M))\qquad\text{(car}
\]
\[
\lambda(\deleted{\unsure{X|x}})(M)\overset{\mathrm{défn}}{=}\psi(M)(x)
=\operatorname{Hom}_{\widehat\Delta}(x,\psi(M))
=\operatorname{Hom}_{\widehat\Delta_{\mathrm{ab}}}(\mathbb Z(x),\psi(M))
\]

donc il doit être représentable, \deleted{\unsure{donc}} (faire $X$
$=\mathbb Z(x)$) donc $\psi\lambda(x)\in A$, donc on est
dans $\underline{\operatorname{Hom}}(\Delta^{\circ},\deleted{\unsure{\mathit{Ab}}}A^{\circ})
\simeq\underline{\operatorname{Hom}}(\Delta,A)$,
on a pratiquement terminé --]

Quel est \inserted{alors} le foncteur commutant aux
$\varinjlim$ part. universel de $\widehat\Delta$ dans cat.
additive et avec $\varinjlim$ ? C'est
\[
X\longrightarrow\mathbb Z(X)=\mathbb Z^{(X)}\quad\text{de }\widehat\Delta\text{ dans }\widehat\Delta_{\mathrm{ab}}
\]
(ou le foncteur ``\uline{abélianisation}'') Sans
doute ça marche pour un topos quelconque.
Mais ici \inserted{(i.e. $\Delta=$ cat. des $\Delta^n$)} cela signifie

\uline{Corollaire} (\unsure{après|à qui|équiv.} : Dold-Puppe \deleted{\unsure{part}}) : \inserted{a)} L'objet
cosimplicial \deleted{$n\mapsto\mathbb Z(\unsure{x})$} $n\longmapsto\mathbb C_{\ast}(\Delta_{\ast}^{n})$
$=\mathbb Z(\Delta_{\ast}^{n})=\mathbb Z^{\Delta_{\ast}^{n}}$ de
$\widehat\Delta_{\mathrm{ab}}\simeq\mathit{Ab}_{\ast}$
est \uline{universel} pour les objets cosimpliciaux
dans des cat. additives avec
$\varinjlim$ quelc, (pour des foncteurs entre
telles cat. qui commutent aux
$\varinjlim$). Donc \deleted{l'objet} \inserted{le complexe des cochaînes} correspondant
